\section{Reference Data Curation}

The current workflow supports two reference sources:

\begin{itemize}
    \item CRAFTED as the primary simulation-reference dataset.
    \item NIST ISODB as an additional experimental or literature reference source.
\end{itemize}

Reference curation is important because not every available isotherm is directly comparable to a simulation result.
Potential problems include different adsorption basis, different temperature, different pressure range, different material naming, duplicated literature entries, and outlier datasets.

\subsection{CRAFTED References}

CRAFTED references are used as structured benchmark references where available.
For the current MOF-5/IRMOF-1 and CO2 workflow, they provide the primary comparison dataset and geometric descriptors such as accessible pore volume.

\subsection{NIST ISODB References}

NIST ISODB entries are searched by material aliases, adsorbate, temperature, and pressure overlap.
The resolver ranks candidate references and records metadata such as DOI, source, adsorption basis, and exclusion reasons.

Candidate references can be excluded from active evaluation if they are detected as outliers or if their adsorption basis is not allowed by the benchmark configuration.
Excluded entries are documented in:

\begin{lstlisting}
evaluations/reference_candidates.json
\end{lstlisting}

\subsection{Adsorption Basis}

The reference parser distinguishes:

\begin{itemize}
    \item \texttt{simulation\_reference},
    \item \texttt{absolute},
    \item \texttt{excess},
    \item \texttt{unknown}.
\end{itemize}

The current evaluation compares references by adsorption basis:

\begin{itemize}
    \item simulation-reference and absolute references are compared to absolute simulation loading,
    \item excess references are compared to excess simulation loading,
    \item unknown-basis references are reported as unknown and compared to absolute loading by default.
\end{itemize}

\subsection{Outlier Handling}

Outliers should not be silently removed.
Instead, the resolver stores excluded references and the reasons for exclusion.
This makes the reference curation reproducible and easier to discuss scientifically.
